% 考古天文学（综述）
% license CCBYNCSA3
% type Wiki

本文根据 CC-BY-SA 协议转载翻译自维基百科\href{https://en.wikipedia.org/wiki/Archaeoastronomy}{相关文章}。


考古天文学（Archaeoastronomy，又作 Archeoastronomy）是一门跨学科（interdisciplinary）[1]或多学科（multidisciplinary）[2]的研究领域，主要探讨古代人类如何“理解天空中的现象、如何利用这些现象，以及天空在其文化中所扮演的角色”[3]。Clive Ruggles 指出，将考古天文学简单地视为“古代天文学的研究”是误导性的，因为现代天文学是一门科学学科，而考古天文学则关注不同文化对天象的象征性与文化性解释[4][5]。考古天文学常与民族天文学（ethnoastronomy）并行研究，后者是一门人类学分支，专门研究当代社会中的观星活动与文化观念。此外，考古天文学还与历史天文学（historical astronomy）}密切相关——  
后者利用古代的天象记录来解决天文学问题，  
以及\textbf{天文学史（history of astronomy）}，  
即通过书面文献来评估古代天文学实践的发展历程[6]。
